% Part: sets-functions-relations
% Chapter: functions
% Section: composition

\documentclass[../../../include/open-logic-section]{subfiles}

\begin{document}

\olfileid[pt]{sfr}{fun}{cmp}
\olsection{Composição de Funções}

\begin{explain}
\oliflabeldef{sfr:fun:inv:sec}{Vimos na \olref[inv]{sec} que a
inversa~$f^{-1}$ de !!a{bijection}~$f$ é ela própria uma função. Outra
operação sobre funções é a composição: p}{P}odemos definir uma nova
função compondo duas funções, $f$ e~$g$, isto é, aplicando primeiro
$f$ e depois~$g$. Claro que isto só é possível se as imagens e os domínios
forem compatíveis, isto é, a imagem de~$f$ tem de ser um subconjunto do
domínio de~$g$. \oliflabeldef{sfr:rel:ops:sec}{Esta operação sobre
funções é o análogo da operação de produto relativo sobre
relações na \olref[rel][ops]{sec}.}{}

Um diagrama pode ajudar a explicar a ideia de composição. Na
\olref{fig:composition}, representamos duas funções $f \colon A \to B$
e $g \colon B \to C$ e a sua composição~$(\comp{f}{g})$. A
função $(\comp{f}{g}) \colon A \to C$ associa cada !!{element} de~$A$
a !!a{element} de~$C$. Especificamos a que !!{element} de~$C$ fica
associado !!a{element} de $A$ do seguinte modo: dada uma entrada $x \in
A$, aplicamos primeiro a função $f$ a~$x$, obtendo algum $f(x)
= y \in B$, depois aplicamos a função $g$ a~$y$, obtendo algum
$g(f(x)) = g(y) = z \in C$.
\begin{figure}[H]
  \olasset[2\olphotowidth]{assets/diagrams/composition.tikz}
  \caption{A composição $g \circ f$ de duas funções $f$ e~$g$.}
  \ollabel{fig:composition}
\end{figure}
\end{explain}

\begin{defn}[Composição]
Sejam $f\colon A \to B$ e $g\colon B \to C$ funções. A
\emph{composição} de $f$ com~$g$ é $\comp{f}{g} \colon A \to C$,
onde $(\comp{f}{g})(x) = g(f(x))$.
\end{defn}

\begin{ex}
Consideremos as funções $f(x) = x + 1$, e $g(x) = 2x$. Como
$(\comp{f}{g})(x) = g(f(x))$, para cada entrada~$x$ deve primeiro tomar-se
o sucessor, depois multiplicar o resultado por dois. Logo a composição
é dada por $(\comp{f}{g})(x) = 2(x+1)$.
\end{ex}

\begin{prob}
Mostre que, se $f \colon A \to B$ e $g \colon B \to C$ são ambas
!!{injective}, então $\comp{f}{g}\colon A \to C$ é !!{injective}.
\end{prob}

\begin{prob}
Mostre que, se $f \colon A \to B$ e $g \colon B \to C$ são ambas
!!{surjective}, então $\comp{f}{g}\colon A \to C$ é !!{surjective}.
\end{prob}

\begin{prob}
Suponhamos $f \colon A \to B$ e $g \colon B \to C$. Mostre que o gráfico
de $\comp{f}{g}$ é $R_f \mid R_g$.
\end{prob}

\end{document}
